
\documentclass[11pt, a4paper]{article}

% Language
\usepackage[english]{babel} 				%option for paper in german
% !TeX spellcheck = en_US

% Font
\usepackage[utf8]{inputenc}					%package to set font encoding
\usepackage[T1]{fontenc}
\usepackage{fontawesome}					
%\usepackage[default]{lato}					%font style
\usepackage{mathpazo}						%font style
\usepackage{eurosym} 						%package to get eurosymbol per \euro
\usepackage{csquotes} 						%flexible quotation marks
\usepackage[super]{nth}						%superscripts for numbers 2nd, 3rd etc
% \usepackage{sansmath}\sansmath				%Turning on sans-serif math mode

% Page Formatting
\setlength{\textheight}{27cm}				%text height
\setlength{\textwidth}{18cm}				%text width
\setlength{\oddsidemargin}{-1cm}			%page margin
\setlength{\evensidemargin}{-1cm} 			%page margin
\setlength{\topmargin}{-3cm}				%top margin
\setlength{\footnotesep}{.5cm}				%space between footnote and text
\usepackage[singlespacing]{setspace} 		%package to adjust line spacing
\setlength{\parindent}{0em}					%set paragraph indentation
\usepackage{pdflscape} 						%package to rotate pages
\usepackage{layouts}						%package to show page formatting settings and convert into different units of measurement

% Bibliography
\usepackage[backend=biber,citestyle=authoryear-comp,bibstyle=authoryear,natbib=true,doi=false,isbn=false,url=false,eprint=false,giveninits=true, maxnames=10, maxbibnames=10,uniquename=false,uniquelist=false]{biblatex}
\addbibresource{~/OneDrive - University of Bristol/10 Literature/Archived/Personal.bib}	%path to library
\AtBeginBibliography{
	\renewcommand*{\mkbibnamefamily}[1]{\textsc{#1}} 			%family names in small caps
}
\renewbibmacro{in:}{}											%No "In: " before journal name
\renewbibmacro*{volume+number+eid}{%							%Volume in brackets
	\printfield{volume}%
	%  \setunit*{\adddot}										% DELETED
	\setunit*{\addnbspace}										% NEW (optional); there's also \addnbthinspace
	\printfield{number}%
	\setunit{\addcomma\space}%
	\printfield{eid}}
\DeclareFieldFormat[article]{number}{\mkbibparens{#1}}


% Hyperlinks
\usepackage{hyperref}
\hypersetup{
	colorlinks,
	linkcolor={red!50!black},
	citecolor={blue!50!black},
	urlcolor={blue!80!black}
}

% Floats
\usepackage{float} 							%allowing for H float specifier
\usepackage{caption}						%define caption font
\captionsetup{labelsep=endash, justification=centering,font=bf,labelfont=sc, skip=5pt, position=top, singlelinecheck=false, parskip=5pt}
\usepackage{placeins}						%package to allow for float barriers
\setlength{\intextsep}{1cm}					%Specify distance of floats to text
\setlength{\floatsep}{1cm} 					%Specify distance between two floats

% Graphics
\usepackage{tikz} 	%package for creating graphics
\usetikzlibrary{calc}


% Tables
\usepackage{booktabs} 						%package for nice tables
\usepackage{multirow} 						%package to combine cells in tables
\usepackage{longtable}						%package to allow for long tables
\usepackage{makecell}						%package to allow line breaks in cells
\usepackage{array}							%package to format tabular environments
%\captionsetup[table]{belowskip=10pt}		%distance text caption
%\captionsetup[longtable]{aboveskip=10pt}	%distance text caption



% Figures
%\captionsetup[figure]{aboveskip=10pt}		%distance text caption
\usepackage{graphicx} 						%package to include graphics
\usepackage{svg}							%include svg images
\usepackage{subfigure}						%package for subfloats

% Math
\usepackage{amsfonts,amssymb, amsthm, mathtools} 	%packages to facilitate writing math
\newtheorem{theorem}{Proposition}					%define "Proposition"-environment
\renewcommand\qedsymbol{$\blacksquare$}				%define qed-symbol
\DeclareMathOperator{\card}{card}
\DeclareMathOperator{\eur}{\text{\euro}}

% Other
\usepackage{comment} 	%package to allow for commenting out entire sections
\usepackage{soul} 		%package to highlight sections
\usepackage{color}		%colorcoding
\newcommand{\hlcyan}[1]{{\sethlcolor{cyan}\hl{#1}}}
\usepackage{enumitem}	%package to allow for enumeration/itemization
\usepackage{titling}	%control package for \maketitle

% JEL/Keywords
\providecommand{\keywords}[1]{\textbf{Keywords:} #1}
\providecommand{\JEL}[1]{\textbf{JEL-Codes:} #1}

%Format Section Headers
\usepackage[explicit]{titlesec} %package for title formatting
\titleformat{\section}{\normalfont\bfseries\large\raggedright\uppercase}{\thesection}{1em}{\large #1}
\titleformat{\subsection}{\normalfont\it\large}{\thesubsection}{1em}{\large #1}
\renewcommand{\thesection}{\arabic{section}}
\renewcommand{\thesubsection}{\thesection.\arabic{subsection}}

%Format Paragraph Headers
%Format Paragraph Headers
\makeatletter
\renewcommand\subparagraph{%
	\@startsection {subparagraph}{5}{\z@ }{0ex \@plus 1ex
		\@minus .2ex}{-1em}{\normalfont \normalsize \it }}%
\makeatother

%Footnotes/Endnotes
\usepackage{endnotes}
%\let\footnote=\endnote

\newcommand{\tabitem}{~~\llap{\textbullet}~~}


\begin{document}
\onehalfspacing
\setlength{\parskip}{.5cm}

%----------------------------------------------------------------------------------------
%	HEADER SECTION
%----------------------------------------------------------------------------------------
\begin{center}
{\LARGE \centering \textbf{Paul Hufe}}\\
{\small \faMousePointer~\href{https://www.paulhufe.net}{www.paulhufe.net}~~~~~~\faTwitter~\href{https://twitter.com/paulhufe?lang=de}{paulhufe}}
\end{center}

\vspace{.5cm}

\section*{Research Statement}

My research interests lie at the intersection of public, labor, and normative economics. Under the overarching theme of equality of opportunity, my research agenda is driven by two main objectives. First, I aim to strengthen the methodological toolkit that is used to quantify the extent of inequality of opportunity in current societies. Thereby, my work connects to the literature branches on intergenerational mobility and inequality measurement. Second, I aim to contribute to our understanding of which circumstantial life factors cause the unequal distribution of life chances. Thereby, my work connects to the literature branches on skill formation and gender inequality.

In the following, I will illustrate three ongoing research projects that relate to this agenda and on which I will work in the coming years.

\paragraph{(Un)Fair inequality in the labor market: A global study}

\subparagraph{Main Collaborators:} \href{http://perseus.iies.su.se/~ialm/}{Ingvild Almås}, IIES Stockholm and \href{https://www.danielweishaar.net/}{Daniel Weishaar}, LMU Munich

\subparagraph{Short Description:} Social justice is one of the dominant issues of our times. According to a recent Gallup poll, 69\% of the world population and 71\% of the British population say that economic differences between rich and poor are unfair in their country. The perceived unfairness of economic outcomes is likely a major reason for increasing social tensions and weakened support for current economic and political institutions. However, despite the increasing focus on fairness, we still lack a proper understanding of perceived unfairness across the world. For example, do citizens perceive inequalities as unfair because actual inequalities do not square with their fairness preferences, or because they have false beliefs about actual inequality and its sources? Which labor market inequalities contribute most to perceived unfairness—e.g., gender gaps, intergenerational disadvantage, or premia for education? Answers to these questions are key to understanding increasing social tensions and equipping policymakers for the implementation of reforms that redress the widespread discontent with current institutions. Yet, due to data limitations current evidence on these questions is scant.

In this project, we will develop a unified framework to consolidate novel data on fairness preferences and beliefs about inequality into measures of perceived unfairness. We will use factorial survey experiments to collect data on fairness preferences and beliefs about inequality from representative samples in 50 countries across the globe covering 67\% of the world's population. In turn, we will match the collected data on beliefs and preferences with corresponding observational data allowing us to compare preferences and beliefs to actual inequalities in these countries. We will use this unique database to answer the following research questions: What drives perceived unfairness in labor markets around the globe---fairness preferences or biased beliefs? How heterogeneous are fairness conceptions and inequality beliefs within countries, and to what extent is this heterogeneity linked to social identity and self-serving biases? To what extent are societal phenomena, such as support for democracy, trust in government institutions, and support for inequality-reducing public policies, linked to perceived unfairness?

\subparagraph{Current Status:} Funded by Riksbankens Jubileumsfond (SEK 7,800,000) and British Academy (GBP 10,000), pending application for UKRI Future Leader Fellowship (GBP 1,500,000), piloting of survey module completed.

\paragraph{The genetic lottery goes to school: Evidence from Norway and the UK}

\subparagraph{Main Collaborators:} \href{https://www.sv.uio.no/psi/english/people/academic/rosacg/index.html}{Rosa Cheesman}, University of Oslo and incoming Ph.D student

\subparagraph{Short Description} Genetic endowments are fixed at conception and they remain constant over the life course of individuals. Therefore, genetic inequality contradicts the widespread goal of providing equal educational opportunities to all members of society. In line with this policy goal, there is a long-standing literature in economics, sociology, and psychology that poses “the question of how well schools reduce the inequity of birth” (Coleman et al., 1966, p.36). This literature has focused on inequalities by socioeconomic status (SES), race, and gender. However, detailed evidence of how well schools address inequality due to genetic endowments is scant.

In this project, we will analyze whether better schools and teaching practices can compensate for genetic disadvantages. To this end, we will combine data from the Norwegian Mother, Father, and Child Cohort Study (MoBa) with Norwegian register data to causally estimate the interaction between genetic endowments and school quality. Since its inception in 1999, MoBa has recruited more than 114,000 women during pregnancy and follows these women and their families over time. Importantly, MoBa avails information on genetic endowments on 44,017 genotyped father-mother-child trios which allow for the identification of causal genetic effects from within-family variation. Furthermore, we will calculate school value-added measures for various educational outcomes from Norwegian register data allowing us to causally estimate school effects. Leveraging the advantages of both data sources, we will provide the first cleanly identified study of gene-environment interactions in the school context.

In a related subproject, we will conduct a similar analysis in the Avon Longitudinal Study of Parents and Children (ALSPAC). ALSPAC has the crucial advantage of containing rich data on students' non-cognitive skills. This data feature allows us to construct measures of teacher value-added that go beyond standard measures of standardized tests for cognitive ability. As a consequence, we can study gene-environment interactions with a focus on the formation of non-cognitive skills---a dimension that has attracted increasing interest in the recent literature on teacher and school effects.  


\subparagraph{Current Status:} Data cleaning completed, Ph.D. student hired as part of the \href{https://euraxess.ec.europa.eu/worldwide/asean/europe-13-doctoral-positions-social-science-genetics-available-across-8}{European Doctoral Network on Social Science Genetics}

\paragraph{Child penalties unpacked: The effect of beliefs and preferences on time allocations in the early stages of parenthood}

\subparagraph{Main Collaborators:} \href{https://sites.google.com/prod/view/sonjasettele/startseite}{Sonja Settele}, CEBI Copenhagen and \href{https://sites.google.com/site/miriamwuestecon/home}{Miriam Wüst}, CEBI Copenhagen 

\subparagraph{Short Description:} In many societies, gender gaps in labor market outcomes have closed significantly in the decades after World War II. However, progress has stalled in recent decades---even in rather gender-equal societies like Denmark. This tendency is commonly attributed to the "child penalty", i.e., the fact that women's labor market outcomes begin to diverge permanently from men's labor market outcomes after the birth of their first child. In fact, current research shows that the "child penalty" is the single most important factor accounting for the remaining gender disparity in labor markets. Yet surprisingly little is known about its fundamental drivers. Since existing evidence suggests that "child penalties" can neither be explained "economically," i.e., by differences in wages between mothers and fathers, nor "biologically," i.e., by the demands of birth and breastfeeding, researchers hypothesize that persistent "gender norms" may be the driving force behind the continued existence of child penalties.

In this project, we will shed light on the fundamental drivers of "child penalties", i.e., the divergence of labor market outcomes among men and women after the birth of their first child. In particular, we will provide systematic evidence on how each of the following factors shapes the gender-specific time allocation of men and women after transitioning into parenthood: (i) parental beliefs about their (future) potential earnings and how these would be affected by taking time off work, (ii) parental beliefs about the productivity of their individual time investments for child well-being, and (iii) parental preference weights for child well-being relative to their own careers. We will collect survey data from a sample of $\approx6,000$ Danish couples at the transition to parenthood. In the survey, we will elicit incentivized measures of beliefs and preferences from mothers and fathers that live in a couple. We will collect two waves of these measures---one during pregnancy around the 12-week pregnancy scan, and one after childbirth around the first birthday of the child. These incentivized survey measures will be matched with the Danish register data allowing us to trace respondents' labor hours, earnings, and parental leave-taking in the years after childbirth.

\subparagraph{Current Status:} Ongoing application for to Rockwool Foundation (DKK 1,100,000), piloting of survey module completed.

\section*{Teaching Statement}

\subparagraph{Teaching Interests:} 
\begin{enumerate}
\item Applied Economics (all levels),
\item Economic Principles (first-year undergraduate),
\item Topics courses on Fairness and Inequality, Gender Economics, Education Economics (third-year undergraduate and above),
\item General interest in interdisciplinary teaching, e.g., the intersection of Economics with Philosophy, Sociology, or Psychology.
\end{enumerate}

\subparagraph{Teaching Philosophy:}
My teaching philosophy rests on three main pillars. 

First, economics suffers from a significant lack of diversity, particularly in terms of the under-representation of women, ethnic minorities, and individuals from disadvantaged socio-economic backgrounds. This lack of diversity is likely to have important repercussions since we risk missing out on valuable contributions and insights that could help address societal challenges. Therefore, I am firmly committed to the adoption of teaching practices that support under-represented groups to view themselves as part of the economics community. For example, to the extent possible, I try to base my teaching on material from a diverse set of authors and choose teaching examples that represent the wider population. Furthermore, I always provide communication channels that allow to preserve anonymity and thereby reduce participation barriers for all students regardless of their background.

Second, I view economics as a social science and also believe that it has to be taught as such. Therefore, I tailor my teaching material to speak to the societal issues that students are passionate about. For example, even in my first-year class on principles of economics, I try to incorporate applied work from the current research frontier to illustrate how fundamental economic concepts can be used to analyze and understand the societal challenges of our times.


Third, I strive to adopt teaching and assessment practices that decompartmentalize the course material and foster comprehensive student learning. For example, I open most of my classes with multiple choice quizzes that do not only cover the material from the previous week but cover all course content to date. This mode of testing nudges students to stay abreast with all course material and to engage in spaced repetitions and interleaved learning practices. Furthermore, I often rely on essay-based assessments that span a wide range of topics and allow me to test students' ability to connect various elements of the course material. 

\end{document}

